\lecture{12}{1. October 2025}{Recap}

\begin{problemwithimage}[0.3\linewidth]{p12_1}{3.11}
  Consider the two-fluid manometer shown.
  Calculate the applied pressure difference.
\end{problemwithimage}

\begin{assumptions}
  \item Incompressible and static liquid.
\end{assumptions}

\begin{derivation}
  To find the pressure difference in the two-fluid manometer we will employ the hydrostatic equation with $z$ being vertically downwards,
  \[ 
  \frac{\mathrm{d}p}{\mathrm{d}z} = \rho g
  .\]
  We let $d$ denote the vertical distance from point 1 to the water level below point 2.
  The pressure contribution from the distance $d$ at either side will cancel out and we therefore need not consider this distance.
  Then, moving from point 1 to point 2 using the hydrostatic equation yields
  \[ 
  p_1 + \rho_{\mathrm{w}} \cdot g \cdot l - \rho_{\mathrm{ct}} \cdot g \cdot l = p_2
.\]
  Solving for the pressure difference $\Delta p = p_2 - p_1$ we get
  \[ 
  \Delta p = \rho_{\mathrm{w}} \cdot g \cdot l - \rho_{\mathrm{ct}} \cdot g \cdot l = \left( \rho_{\mathrm{w}} - \rho_{\mathrm{ct}} \right) gl
  .\]
  Inputting the different quantities yields
  \[ 
  \Delta p = \left( \qty{1000}{\kg\per\cubic\m} - \qty{1590}{\kg\per\cubic\m} \right) \cdot \qty{9,81}{\m\per\s\squared} \cdot \qty{10,2e-3}{\m} = \qty{-59,1}{\Pa} 
  .\]
  Therefore the pressure at point 2 will be \qty{59,1}{\Pa} lower than that at point 1.
\end{derivation}

\finalresult{\Delta p = \qty{-59,1}{\Pa} }


\begin{problemwithimage}{p12_2}{4.10}
  Find $V$ for this mushroom cap on a pipeline.
\end{problemwithimage}

\begin{assumptions}
  \item Steady and incompressible flow
\end{assumptions}

\begin{derivation}
  The continuity equation for steady flow is
  \[ 
    0 = \int_{\mathrm{CS}} \rho \Vec{V} \cdot \, \mathrm{d} \Vec{A}
  .\]
  Since the density is constant the flow rate $Q = V \cdot A$ is constant.
  Hence,
  \[
    Q_{\mathrm{in}} = Q_{\mathrm{out}}
  .\]
  The flow rate out of the cap can also be written as
  \[ 
  Q_{\mathrm{out}} = V \cdot \cos \ang{45} \cdot A_{\mathrm{out}}
.\]
  Combining these yields
  \begin{align*}
    Q_{\mathrm{in}} &= V \cdot \cos \ang{45} \cdot A_{\mathrm{out}} \\
    V &= \frac{Q_{\mathrm{in}}}{\cos \ang{45} \cdot A_{\mathrm{out}}} \\
      &= \frac{\qty{3}{\cubic\m\per\s}}{\cos \ang{45} \cdot \pi \cdot \left( \left( \qty{2}{\m}  \right)^2 - \left( \qty{1,8}{\m}  \right)^2 \right)} \\
      &= \qty{1,78}{\m\per\s} 
  .\end{align*}
\end{derivation}

\finalresult{V = \qty{1,78}{\m\per\s} }

\begin{problemwithimage}{p12_3}{4.4}
  Fluid with \qty{1040}{\kg\per\cubic\m} density is flowing steadily through the rectangular box shown.
  Given $A_1 = \qty{0,046}{\m\squared}$, $A_2 = \qty{0,009}{\square\m}$, $A_3 = \qty{0,056}{\square\m}$, $\Vec{V}_1 = \left( \qty{3}{\m\per\s}  \right)\hat{i} $, and $\Vec{V}_2 = \left( \qty{6}{\m\per\s}  \right) \hat{j}$. Determine the velocity $\Vec{V}_3$.
\end{problemwithimage}

\begin{assumptions}
  \item Incompressible, steady and uniform flow.
\end{assumptions}

\begin{derivation}
  Due to the assumptions of uniform and incompressible steady flow the continuity equation reduces to
  \[ 
    \sum_{CS} \Vec{V} \cdot \Vec{A} = 0 
  .\]
  Using the convention of $\Vec{A}$ being outwards we get
  \begin{align*}
    V_1 \cdot A_1 + V_2 \cdot A_2 + V_3 \cdot A_3 &= 0 \\
    V_3 &= -\frac{V_1 \cdot A_1 + V_2 \cdot A_2}{A_3} \\
              &= \frac{\qty{3}{\m\per\s} \cdot \qty{0,046}{\square\m} - \qty{6}{\m\per\s} \cdot \qty{0,009}{\square\m}}{\qty{0,056}{\square\m} } \\
              &= \qty{1,5}{\m\per\s} 
  .\end{align*}
  Using trigonometry this can be written on Cartesian form as
  \begin{align*}
    V_{3x} &= V_3 \cdot \cos \ang{-30} = \qty{1,5}{\m\per\s} \cdot \cos \ang{-30} = \qty{1,3}{\m\per\s} \\
    V_{3y} &= V_3 \cdot \sin \ang{-30} = \qty{1,5}{\m\per\s} \cdot \sin \ang{-30} = \qty{-0,75}{\m\per\s}  \\
    \Vec{V}_3 &= V_{3x} \hat{i} + V_{3y} \hat{j} = \qty{1,3}{\m\per\s} \, \hat{i} - \qty{0,75}{\m\per\s} \, \hat{j}
  .\end{align*}
\end{derivation}

\finalresult{\Vec{V}_3 = \qty{1,3}{\m\per\s} \hat{i} - \qty{0,75}{\m\per\s} \hat{j}}


\begin{problem}{5.19}
  A velocity field is given by $\Vec{V} = 2 \hat{i} - 4 x \hat{j} \, \unit{\m\per\s}$.
  Determine an equation for the streamline.
  Calculate the vorticity of the flow.
\end{problem}

\begin{assumptions}
  \item The flow is steady and incompressible
\end{assumptions}

\begin{derivation}
  The stream function for a steady velocity field $(u,v)$ is defined as
  \[ 
  u \equiv \frac{\partial \psi}{\partial y} \qquad v \equiv - \frac{\partial \psi}{\partial x} 
  .\]
  Using $u = \qty{2}{\m\per\s}$ we get
  \[ 
  \frac{\partial \psi}{\partial y} = 2
  .\]
  By integration we obtain
  \[ 
  \psi = 2y + f(x)
  .\]
  Using $v = - 4 x$ we get
  \[ 
  - \frac{\partial \psi}{\partial x} = -4x \implies \frac{\partial \psi}{\partial x} = 4x
  .\]
  By integration this becomes
  \[ 
  \psi = 2 x^2 + f(y)
  .\]
  Hence, the stream function is:
  \[ 
  \psi = 2x^2 + 2y + c
  \]
  where $c$ is a constant. 
  The vorticity $\xi$ is given by
  \begin{align*}
    \xi &\equiv \nabla \times \Vec{V} \\
        &= \left( \frac{\partial v}{\partial x} - \frac{\partial u}{\partial y}  \right) \, \hat{k} \\
        &= \left( -4 - 0 \right) \, \hat{k} \\
        &= -4 \, \hat{k}
  .\end{align*} 
\end{derivation}

\finalresult{
  \begin{aligned}
    \psi(x,y) &= 2x^2 + 2y + c \\
    \Vec{\xi} &= - 4 \, \hat{k}
  \end{aligned}
}

